\documentclass[11pt,a4paper]{article}
\usepackage[utf8]{inputenc}
\usepackage{amsmath,amssymb,amsfonts}
\usepackage{geometry}
\usepackage{hyperref}

\geometry{margin=1in}

\title{\textbf{\Huge Equal Temperament Pitch \& Frequency Conversions}\\[0.5em]
\Large Mathematical Models of \texttt{[mtof]} and \texttt{[ftom]} Nodes}
\author{\textbf{Time Dilation DAW (v0.0.1) Engineering Group}}
\date{\today}

\begin{document}

\maketitle

\begin{abstract}
This document formulates the logarithmic and exponential relationships between MIDI Note numbers $m \in [0, 127]$, fundamental frequencies $f \in \mathbb{R}^+$ (in Hz), microtonal cent tuning ratios, and Doppler frequency shifts under relativistic coordinate time dilation.
\end{abstract}

\section{MIDI Note to Frequency Conversion (\texttt{[mtof]})}
In 12-tone equal temperament (12-TET) with concert pitch reference $A_4 = 440\,\text{Hz}$ at MIDI Note $m = 69$, the fundamental frequency $f(m)$ is defined exponentially:
\begin{equation}
f(m) = 440 \cdot 2^{\frac{m - 69}{12}}
\end{equation}
For an arbitrary reference concert pitch $A_4 = f_0$:
\begin{equation}
f(m, f_0) = f_0 \cdot 2^{\frac{m - 69}{12}}
\end{equation}

\section{Frequency to MIDI Note Conversion (\texttt{[ftom]})}
The inverse mapping from fundamental frequency $f > 0$ to MIDI Note number $m(f)$ is derived by taking the base-2 logarithm:
\begin{equation}
m(f) = 69 + 12 \cdot \log_2\left(\frac{f}{440}\right)
\end{equation}
Using standard natural logarithms:
\begin{equation}
m(f) = 69 + \frac{12}{\ln 2} \cdot \ln\left(\frac{f}{440}\right)
\end{equation}

\section{Cent Tuning \& Transposition Ratios}
The frequency multiplier ratio $R(c)$ for a microtonal interval of $c$ cents ($100\,\text{cents} = 1\,\text{semitone}$) is:
\begin{equation}
R(c) = 2^{\frac{c}{1200}}
\end{equation}
Given a base frequency $f_{\text{base}}$ and semitone shift $\Delta m$ with fine detune cents $c$, the modulated frequency is:
\begin{equation}
f_{\text{final}} = f_{\text{base}} \cdot 2^{\frac{\Delta m + c/100}{12}}
\end{equation}

\section{Relativistic Doppler Pitch Shift under Time Warping}
When an audio node is driven by a time dilation factor $\gamma(t)$, the apparent perceived playback frequency $f_{\text{perceived}}$ experiences an acoustic Doppler shift:
\begin{equation}
f_{\text{perceived}}(t) = f_{\text{base}} \cdot |\gamma(t)|
\end{equation}
Under retrograde time flow ($\gamma(t) < 0$), the frequency magnitude remains $|f_{\text{perceived}}| = f_{\text{base}} \cdot |\gamma(t)|$ while the phase trajectories invert direction.

\end{document}
